\documentclass[12pt]{article}
\usepackage[T1]{fontenc}
\usepackage[utf8]{inputenc}
\usepackage{lmodern}
\usepackage{textcomp}
\usepackage{enumitem}
\usepackage{geometry}
\geometry{margin=1in}
\begin{document}
\begin{center}
Answer Key - Psychology - Graduate Level Assessment 18 \
\small (Answers are ordered exactly as in the question paper.)
\end{center}

\section*{Multiple Choice}
\begin{enumerate}[label=\arabic*.]
\item \textbf{Answer:} A
\\ \textit{Options:} A) Random assignment, temporal precedence; B) Nonspuriousness, temporal precedence; C) Generalizability, causality; D) Temporal precedence, nonspuriousness; E) Temporal precedence, random selection
\\ \textit{Note:} Random assignment helps ensure that the groups in an experiment are equivalent at the start, enhancing internal validity by controlling for confounding variables, while temporal precedence establishes a cause-and-effect relationship by demonstrating that the cause precedes the effect, which is essential for external validity as it allows the findings to be generalized to other contexts.
\item \textbf{Answer:} A
\\ \textit{Options:} A) are definitional rather than conceptual; B) check memorization rather than critical thinking; C) encourage rote learning rather than comprehension; D) require less time to answer; E) enhance retention of information
\\ \textit{Note:} Multiple-choice questions tend to focus on recalling specific definitions or facts, making them less cognitively demanding compared to fill-in or completion questions that require a deeper understanding of concepts and the ability to generate responses independently.
\item \textbf{Answer:} C
\\ \textit{Options:} A) Stronger infant immunity solely based on the duration of breastfeeding; B) specific method used in feeding; C) the attitude the parents have towards the child during feeding time; D) Predictable sleep patterns established immediately through breastfeeding; E) Enhanced cognitive development exclusively due to nutrient composition
\\ \textit{Note:} Research indicates that the emotional and psychological interactions during breastfeeding significantly influence the parent-infant bond and the child's emotional development, making parental attitudes during feeding a critical factor.
\item \textbf{Answer:} C
\\ \textit{Options:} A) blue; B) red; C) dark; D) white; E) shiny
\\ \textit{Note:} Objects that absorb light appear dark because they do not reflect sufficient light to be perceived as bright or colored by the human eye, leading to a perception of reduced brightness.
\item \textbf{Answer:} C
\\ \textit{Options:} A) adjustment; B) homeostasis; C) exhaustion; D) stressor detection; E) resolution
\\ \textit{Note:} The exhaustion stage of Selye's General Adaptation Syndrome (GAS) theory is when the body's resources are depleted after prolonged stress, leading to decreased immune function and increased susceptibility to diseases.
\end{enumerate}

\section*{True / False}
\begin{enumerate}[label=\arabic*.]
\setcounter{enumi}{5}
\item \textbf{Answer:} True
\\ \textit{Options:} A) True; B) False
\\ \textit{Note:} Research indicates that young Caucasian women often experience higher rates of mental health issues and suicidal ideation, leading them to seek out and utilize suicide prevention services more frequently than other demographic groups.
\item \textbf{Answer:} True
\\ \textit{Options:} A) True; B) False
\\ \textit{Note:} The Kuder Occupational Interest Survey is developed from a different theoretical framework that emphasizes specific interest patterns rather than being directly derived from Holland's theory of vocational interests.
\item \textbf{Answer:} False
\\ \textit{Options:} A) True; B) False
\\ \textit{Note:} Freud's theory suggests that the id operates on the pleasure principle and is often counterbalanced by the ego and superego, meaning that Andy's behavior reflects a more complex interplay of these structures rather than a straightforward dominance of the id.
\item \textbf{Answer:} True
\\ \textit{Options:} A) True; B) False
\\ \textit{Note:} The 'miracle question' is an initial intervention in solution-focused therapy designed to help clients envision their desired future and identify the steps needed to achieve their goals, making it a fundamental technique in this therapeutic approach.
\item \textbf{Answer:} False
\\ \textit{Options:} A) True; B) False
\\ \textit{Note:} Donald Super's life cycle wheel incorporates various life roles, such as family, community, and leisure, alongside the sequential phases of career development, highlighting the interplay between career and personal life.
\end{enumerate}

\section*{Long Form Response}
\begin{enumerate}[label=\arabic*.]
\setcounter{enumi}{10}
\item \textbf{Answer:} Individuals with high internality on the Internal-External Control Scale exhibit a strong belief in their ability to influence outcomes through their own actions and decisions. This characteristic is closely linked to a heightened achievement orientation, as such individuals tend to set ambitious goals and persistently pursue them, believing that success is a direct result of their efforts. Additionally, they often demonstrate an independent decision-making style, relying on their own judgment rather than succumbing to external pressures or influences. This independence fosters resilience and adaptability, enabling them to navigate challenges effectively. Consequently, individuals with high internality are more likely to engage in tasks that require skill and effort, rather than those left to chance, ultimately leading to greater success in various domains of life.
\\ \textit{Note:} correct because it accurately connects high internality on the Internal-External Control Scale to a strong belief in personal agency, which fosters a proactive achievement orientation and an independent decision-making style.
\item \textbf{Answer:} Traditional research typically focuses on generating knowledge through hypothesis testing, often utilizing controlled experimental designs and statistical analysis to draw conclusions about specific variables. In contrast, program evaluation employs unique methodologies that blend both quantitative and qualitative approaches, such as surveys, interviews, focus groups, and observational studies, to assess the effectiveness and impact of specific programs or interventions. This distinction is crucial in psychology because program evaluation not only seeks to determine whether a program works but also explores the context, implementation processes, and participant experiences, providing a comprehensive understanding of the program's effects. By integrating diverse data collection methods, program evaluation offers insights that can inform practice and policy, ultimately enhancing the relevance and application of psychological research in real-world settings.
\\ \textit{Note:} The answer correctly identifies that traditional research emphasizes hypothesis testing and controlled designs, while program evaluation uses diverse methodologies to assess program effectiveness, highlighting the importance of these distinctions in understanding how psychological interventions are evaluated and improved.
\end{enumerate}

\vfill
\noindent\textit{End of Answer Key}
\end{document}